% =========================================
% CHAPTER 1: INTRODUCTION
% =========================================
\chapter{Introduction}
\label{ch:intro}

\section{The Problem}

Bitcoin~\cite{nakamoto2008bitcoin} solved a specific problem: two strangers can exchange value without a bank. Ethereum~\cite{wood2014ethereum} generalized that to programmable computation. Both networks hit the same wall. Bitcoin processes 7 transactions per second. Ethereum processes 15. Visa processes thousands per second. This gap matters because it means blockchain fees spike under demand and most applications cannot use the network reliably.

The reason is simple: every node must validate every transaction. This keeps the network secure and decentralized, but it is slow. You cannot make it faster without breaking something.

Layer~2 systems try a different approach. Instead of validating every transaction, Layer~1 validates a summary—a single commitment representing thousands of transactions batched together. The commitment is how Layer~2 anchors security to Layer~1. If the commitment is wrong, Layer~1 rejects it. If it is right, the transactions are secure.

Two approaches emerged. ZK~Rollups compute a cryptographic proof showing the commitment is valid, then submit the proof to Layer~1. Layer~1 checks the proof and accepts the commitment immediately. Optimistic~Rollups submit the commitment without proof. They give anyone a window to challenge it if it is wrong. If no valid challenge appears, the commitment becomes final.

Which is better? It depends. ZK Rollups finalize faster but proof computation is expensive. Optimistic Rollups are cheaper but users wait days for settlement. We lack empirical data on how they actually compare under stress.

\section{What Is Missing}

Academic papers describe the theory. Production systems achieve throughput at scale. But nobody measures what happens when things break.

When Layer~1 undergoes a chain reorganization (it happens), which finality guarantee is more robust? When a sequencer crashes or disappears, what can users do and what does recovery cost? When a commitment contains an error, how does each system detect and recover?

These are not hypothetical. Chain reorganizations have occurred. Sequencer outages have happened. Yet the literature offers few concrete answers on how ZK versus Optimistic compare under stress.

Most evaluation papers measure throughput when the network is uncongested, latency when sequencers work correctly, and costs under ideal conditions. It is useful, but it is incomplete.

\section{What This Thesis Does}

This thesis builds a testing framework and uses it to measure both rollup types under normal operation and deliberate failures. The goal is to characterize when each approach is preferable.

The framework runs both sequencers against the same Layer~1 node and measures what happens. We simulate sequencer crashes. We inject Layer~1 delays. We submit invalid states. For each scenario, we record finality latency, recovery time, gas cost, and whether data is lost.

Everything is open-source and deployable with a single Docker command. This means other researchers can verify the results, extend the evaluation, and build on the work.

\section{Research Questions}

The evaluation answers four specific questions.

First, how long does a transaction actually take to become irreversible? Theory says ZK Rollups finalize in seconds and Optimistic Rollups finalize in days. What do we measure in practice?

Second, when a sequencer crashes, what is the cost to recover? Do users lose their transactions? How much gas is required to get the system working again?

Third, what are the actual trade-offs? ZK Rollups compute proofs. Optimistic Rollups wait for challenge windows. Where do the real costs come from?

Fourth, can we build a fair comparison on standard developer hardware? If yes, the framework is useful. If no, it is a toy.

\section{Why a Framework?}

Why not deploy both systems on mainnet and compare? Because environmental noise masks architectural differences. Gas prices change. Network conditions vary. Sequencer implementations differ in hundreds of ways. A controlled framework isolates the essential differences.

Why not just do formal analysis? Because real systems diverge from theory. Proof generation has overhead that depends on batch size, hardware, and implementation choices. Fraud proofs require network propagation delays that formal analysis does not capture. Real bugs occur in real systems. Formal analysis gives bounds; measurement gives data.

\section{Contributions}

This work contributes three things. An open-source testing framework with configurable parameters, allowing researchers to test different batch sizes, polling frequencies, and gas allocations. Empirical data on finality, recovery cost, and resource overhead for both architectures under controlled conditions. A documented interface specification so researchers can implement custom sequencers and test them against the framework contracts.

\section{Thesis Structure}

Chapter~\ref{ch:background} covers the technical foundations: how blockchains work, why they are slow, and how ZK and Optimistic security models differ.

Chapter~\ref{ch:relatedwork} surveys existing research and identifies the specific gaps this work addresses.

Chapter~\ref{ch:framework} describes the framework architecture and design choices.

Chapter~\ref{ch:implementation} details the code: the contracts, sequencer processes, and metrics collection.

Chapter~\ref{ch:evaluation} reports the empirical results.

Chapter~\ref{ch:conclusion} synthesizes the findings and discusses implications.
